\chapter*{À propos du projet Open Logic}
\addcontentsline{toc}{chapter}{À propos du projet Open Logic}

Le \textit{Open Logic Text} est un manuel libre et collaboratif de
métalogique et de méthodes formelles. Il commence à un niveau
intermédiaire, après un cours d'introduction à la logique formelle.
Il s'adresse à des lecteurs sans formation spécialisée en
mathématiques, notamment aux étudiants en philosophie et en
informatique. L'exposé reste rigoureux.

La couverture et le niveau de détail varient selon les sujets.
Si vous trouvez une erreur ou avez une suggestion,
\href{https://github.com/OpenLogicProject/OpenLogic/wiki/Contributing}{signalez-la
à l'équipe du projet}.

Le projet fonctionne dans l'esprit des logiciels libres. Le texte
est accessible gratuitement et ses sources LaTeX sont diffusées sous
licence Creative Commons Attribution. Cette licence autorise chacun à
télécharger, utiliser, modifier, réorganiser, convertir et redistribuer
l'ouvrage, à condition d'en attribuer correctement la paternité.
Pour plus de renseignements, consultez le site du projet Open Logic :
\href{http://openlogicproject.org/}{openlogicproject.org}.
